\section{はじめに}
画像中の物体に対するアノテーションは，物体の位置を表すバウンディングボックスと，物体が何であるかを表すクラスの組から構成される．
YOLO\cite{ref_YOLO}やDetection Transformer（DETR）\cite{ref_DETR}に代表される物体認識モデルは，あらかじめ定めたクラスの集合で学習されるため，学習時のクラスの範囲では高い精度で物体を認識できる．
一方，学習時に定めたクラスより細かい粒度，例えば「dog」に対する犬種のようなクラスを付けるためには，そのクラスを含むデータでモデルを学習し直す必要がある．

画像と文章を同じ空間に埋め込むContrastive Language-Image Pre-training（CLIP）\cite{ref_CLIP}やSigLIP\cite{ref_SigLIP}のような文章・画像埋め込みモデルを用いると，任意の単語と画像の類似度を追加の学習なしに計算できる．
また，語彙の上位語と下位語の関係を記述したWordNet\cite{ref_WordNet}を用いると，ある単語より細かい粒度の単語の集合を得ることができる．
OWL-ViT\cite{ref_OWLViT}やOWLv2\cite{ref_OWLv2}のようなオープンボキャブラリの物体認識モデルは，これらの埋め込みを用いて，テキストで与えたクラスの物体を検出できる．
しかし，通常はアノテーションに用いるクラスの集合を利用者が事前に与える必要があり，学習時に無い物体を数千語の語彙から1語で命名すると，誤りが多い．
例えば4.4節で示すように，LVIS\cite{ref_LVIS}の人手の注釈を学習に用いていないOWLv2では，COCO外の物体に対する7806語からの1位の単語の正解率は18.0\%であった．
一方，正しい単語が1位でなくても，その上位語（例えばparkaに対するjacketやclothing）は正しいことが多い．
したがって，学習時に無い物体に対しては，常に細かい1語を出力するのではなく，正しい可能性が高い最も細かい語を選び，その正しさの確率を付けることが望ましい．

本稿では，物体認識モデルとオープンボキャブラリの物体認識モデルの領域の特徴，およびWordNetを組み合わせ，モデルの再学習を行わずに，学習時に無い物体を検出し，粗い語から細かい語までの粒度で命名する粒度可変アノテーションを提案する．
提案手法は，領域の特徴から求めた語彙全体の確率をWordNetの上位語へ集約し，集約した確率が閾値以上の間だけ上位語から下位語へ進むことで，確信度に応じた粒度の単語を選ぶ．
さらに，各アノテーションと，粗い語から選んだ単語までの経路上の各単語に，正しさの確率を表すスコアを付ける．
これにより，利用者はスコアの閾値によって，出力の数と粒度を必要な正解率に応じて選ぶことができる．
スコアの重みは検証用の画像のみで求める14個であり，物体認識モデルと埋め込みモデルは再学習しない．

本稿の貢献は以下の3点である．
第一に，再学習なしで学習時に無い物体を検出し，確信度に応じた粒度で命名する手法と，そのアノテーションおよび経路上の各単語の正しさの確率を求める方法を定式化する．
第二に，WordNetのSynsetが付与されたLVISのアノテーションと，評価に用いるカテゴリを学習時に見ていない重みを用いて，COCO外の物体に対する性能を評価する．
その際，階層を用いずに1語を選ぶ方法や，評価の語彙を事前に与える方法と比較し，階層的な指標とラベルの正誤のみによる指標の両方を，ブートストラップによる信頼区間とともに示す．
第三に，誤りを種類ごとに分解し，提案手法の精度を制限する要因が，物体でない領域や物体の一部への命名にあることを示す．
また，予備実験として，COCOの80クラスの範囲で切り抜き画像と文章・画像埋め込みモデルによるクラスの推定を評価し，学習時のクラスの範囲では物体認識モデル自身のクラスを用いるべきであることを示す．

本稿の構成は以下のとおりである．
2章で関連研究，3章で提案手法を述べる．
4章では，予備実験と，COCO外の物体に対する評価の結果を示す．
5章で考察を行い，6章で本稿をまとめる．
